\section{Kirchhoff--Love 薄板}
\label{chap:thin-plate-theory}

弦和膜的弹性能依赖位移的一阶空间导数，变分后得到二阶波动方程；薄板弯曲能却由曲率控制，曲率已经含有二阶导数，因此运动方程会再多两个空间导数。要从三维弹性得到这个四阶方程，关键不是背诵 Kirchhoff--Love 假设，而是判断厚度 $h$ 与面内尺度 $L$ 的比值怎样压低横向剪切和厚度伸缩。

因此先保留完整的三维位移场，用 $h/L\ll1$ 给各应变分量分阶，再沿厚度积分。二维弯曲能、弯矩、Euler--Lagrange 方程和自然边界项都会从同一次变分中出现。得到这个二维边值问题以后，下一节再把圆板和长方形的几何条件代入，求出共振频率与节线。
\subsection{从三维线弹性到 Kirchhoff--Love 薄板}

三维小应变弹性的平衡方程是二阶的；把厚度方向消去以后，弯曲能却含有中面位移的二阶导数，Euler--Lagrange 方程因而升为四阶。本节从这一步升阶出发，并把每一次降维都对应到一个明确的小参数。中面区域记为 \(\mathcal A\subset\mathbb R^2\)，厚度坐标为 \(-h/2<z<h/2\)，典型面内变化尺度为 \(L\)。记 \(\vb r=(x_1,x_2)\)、\(\nabla_t=(\partial_1,\partial_2)\)，\(\dd a\) 为中面面积元。薄参数定义为
\begin{equation}
 \epsilon_h:=\frac{h}{L}\ll1.
 \label{eq:phys-cont-kirchhoff-slenderness}
\end{equation}
三维位移为 \(\vb u(\vb r,z,t)\)，小应变为
\[
 \varepsilon_{ij}=\frac12(\partial_i u_j+\partial_j u_i).
\]
先不假定各向同性。令 \(\mathsf C^{3\mathrm D}(\vb r,z)\) 为三维线弹性张量、\(\rho_{\rm m}(\vb r,z)\) 为质量密度、\(\vb b(\vb r,z,t)\) 为单位体积体力密度，则在体力项显式写出、表面牵引由边界虚功处理时，线性三维作用量为
\begin{equation}
 \mathcal S_{3\mathrm D}
 =\int\dd t\int_{\mathcal A}\dd a\int_{-h/2}^{h/2}\dd z
 \left[
 \frac{\rho_{\rm m}}{2}|\partial_t\vb u|^2
 -\frac12\vb{\varepsilon}:\mathsf C^{3\mathrm D}:\vb{\varepsilon}
 +\vb b\cdot\vb u
 \right].
 \label{eq:phys-cont-3d-elastic-action}
\end{equation}
Kirchhoff--Love 理论不是另一个独立模型，而是\eqnrefs{eq:phys-cont-3d-elastic-action} 在弯曲、长波、薄厚度极限中的首阶二维有效作用量。三维 Cauchy 应力记为 \(T_{ij}:=C^{3\mathrm D}_{ijkl}\varepsilon_{kl}\)，因此后文的 \(T_{\alpha\beta},T_{\alpha z},T_{zz}\) 都是同一应力张量的分量。

取横向挠度尺度 \(w_\star\)，并针对弯曲支令
\[
 u_z\sim w_\star,
 \qquad
 u_\alpha\sim \frac{h}{L}w_\star.
\]
为做量纲估计，记 \(E_\star\) 为三维刚度张量在相关应变方向上的代表性弹性模量、\(\rho_{{\rm m},\star}\) 为代表性质量密度；二者都假定在 \(\epsilon_h\to0\) 时保持有限且非零。它们只用于尺度计数，不预设各向同性；真正的 Young 模量 \(E\) 只在后面的各向同性特例中出现。那么弯曲产生的面内应变尺度为
\[
 \varepsilon_{\alpha\beta}^{(b)}
 \sim \frac{h w_\star}{L^2},
\]
而若横向剪切不发生抵消，
\[
 \varepsilon_{\alpha z}
 =\frac12(\partial_\alpha u_z+\partial_z u_\alpha)
 \sim \frac{w_\star}{L}.
\]
因此在同一面积 \(L^2\) 上，剪切能与弯曲能的尺度分别为
\[
 U_{\rm shear}\sim E_\star h w_\star^2,
 \qquad
 U_{\rm bend}\sim E_\star\frac{h^3w_\star^2}{L^2},
\]
从而
\begin{equation}
 \frac{U_{\rm shear}}{U_{\rm bend}}
 \sim\epsilon_h^{-2}.
 \label{eq:phys-cont-kl-shear-scaling}
\end{equation}
\eqnrefs{eq:phys-cont-kl-shear-scaling} 首先只排除了未抵消的 \(O(w_\star/L)\) 横向剪切：若归一化弯曲能在 \(\epsilon_h\to0\) 时保持有限，就必须有
\[
 \partial_z u_\alpha+\partial_\alpha u_z
 =O(\epsilon_h)\frac{w_\star}{L}.
\]
仅凭能量阶数，这个估计还允许残余剪切与弯曲应变同阶，因此还不能把它直接等同于 Kirchhoff--Love 约束。对厚度方向是材料反射对称轴、弹性模量不随 \(h\) 奇异缩放的板，在距边缘 \(O(h)\) 之外再使用三维平衡方程和上下表面的无牵引条件可以把阶数进一步收紧。弯曲主应力量级为
\[
 T_{\alpha\beta}\sim E_\star\frac{hw_\star}{L^2},
\]
而
\[
 \partial_\beta T_{\alpha\beta}+\partial_zT_{\alpha z}
 =\rho_{\rm m}\,\partial_t^2u_\alpha
\]
在弯曲时间尺度上给出
\[
 T_{\alpha z}=O(\epsilon_h)T_{\alpha\beta},
 \qquad
 T_{zz}=O(\epsilon_h^2)T_{\alpha\beta}.
\]
因此横向剪切应变是 \(O(\epsilon_h)\) 倍的弯曲应变，其能量相对弯曲能只占 \(O(\epsilon_h^2)\)。这一步才允许在首阶外层理论中令 \(\varepsilon_{\alpha z}=0\)，同时保留由零横向法向应力诱导的 \(\varepsilon_{zz}\)。相应运动学可写为
\begin{equation}
 \vb u_t(\vb r,z,t)
 =\vb u_t^{(0)}(\vb r,t)-z\nabla_t w(\vb r,t)
 +O(\epsilon_h^3w_\star),
 \qquad
 u_z=w+O(\epsilon_h^2w_\star),
 \label{eq:phys-cont-kl-kinematics}
\end{equation}
其中误差阶数指远离三维边界层的弯曲外层展开。这就是 Kirchhoff--Love 运动学：截面直线在首阶保持直且垂直于中面。此时不能立刻把中面面内位移删掉，因为一般非对称层合材料会把膜应变与曲率耦合。先定义中面应变与曲率
\begin{equation}
 e^{(0)}_{\alpha\beta}
 :=\frac12(\partial_\alpha u^{(0)}_\beta+\partial_\beta u^{(0)}_\alpha),
 \qquad
 \kappa_{\alpha\beta}:=-\partial_\alpha\partial_\beta w,
 \label{eq:phys-cont-kl-strain}
\end{equation}
则首阶面内应变具有仿射厚度结构
\[
 \varepsilon_{\alpha\beta}(z)
 =e^{(0)}_{\alpha\beta}+z\kappa_{\alpha\beta}.
\]
因此 Kirchhoff--Love 约束本身只消除了首阶横向剪切；是否能够进一步只研究纯弯曲，必须由二维构成张量的膜--弯曲耦合决定。

还需区分“\(u_z\) 首阶不随 \(z\) 变化”和“强行令 \(\varepsilon_{zz}=0\)”。后者一般错误：自由表面要求的是横向法向应力松弛。前面已经由平衡方程把横向剪切降到高阶；若厚度方向还是材料反射对称轴，则含奇数个 \(z\) 指标的弹性耦合分量消失，横向剪切不会在首阶与面内应变重新耦合。此时只需对 \(\varepsilon_{zz}\) 做平面应力消元。把三维能量中仍处于首阶的部分写成
\[
 W_{3\mathrm D}^{(0)}
 =\frac12 C_{\alpha\beta\gamma\delta}
 e_{\alpha\beta}e_{\gamma\delta}
 +C_{\alpha\beta zz}e_{\alpha\beta}\varepsilon_{zz}
 +\frac12C_{zzzz}\varepsilon_{zz}^2,
 \qquad e_{\alpha\beta}:=\varepsilon_{\alpha\beta}.
\]
在给定面内应变下令 \(T_{zz}=\partial W_{3\mathrm D}^{(0)}/\partial\varepsilon_{zz}=0\)，得到
\begin{equation}
 \varepsilon_{zz}^{\star}
 =-\frac{C_{zz\alpha\beta}}{C_{zzzz}}e_{\alpha\beta},
 \qquad
 Q_{\alpha\beta\gamma\delta}
 :=C_{\alpha\beta\gamma\delta}
 -\frac{C_{\alpha\beta zz}C_{zz\gamma\delta}}{C_{zzzz}}.
 \label{eq:phys-cont-plane-stress-schur}
\end{equation}
\(\mathsf Q\) 是平面应力约化刚度；它是三维构成张量关于横向法向应变的 Schur 补，而不是额外假设。若三维弹性能对非零应变严格正定，则 \(C_{zzzz}>0\)，相应 Schur 补在面内应变子空间上同样正定，因此后面的膜刚度与弯曲刚度继承稳定性。若材料没有上述厚度反射对称性，面内--横向剪切耦合可能在同一阶出现，此时必须使用更一般的各向异性板理论，不能把下面的 Kirchhoff--Love 构成式机械套用。

对允许沿中面按宏观尺度 \(L\) 非均匀、并可随厚度变化但逐点满足上述对称性的材料，把平面应力刚度的前三个厚度矩定义为
\begin{equation}
 \mathsf A(\vb r):=\int_{-h/2}^{h/2}\mathsf Q(\vb r,z)\,\dd z,
 \qquad
 \mathsf B(\vb r):=\int_{-h/2}^{h/2}z\mathsf Q(\vb r,z)\,\dd z,
 \qquad
 \mathsf D(\vb r):=\int_{-h/2}^{h/2}z^2\mathsf Q(\vb r,z)\,\dd z.
 \label{eq:phys-cont-general-plate-rigidity}
\end{equation}
它们分别是膜刚度、膜--弯曲耦合刚度和弯曲刚度；在 SI 制下其分量单位依次为 \(\si{\newton\per\metre}\)、\(\si{\newton}\)、\(\si{\newton\metre}\)。把 \(\vb e^{(0)}\) 与 \(\vb{\kappa}\) 作为独立二维应变变量，单位中面面积的二维能量密度为
\begin{equation}
 \mathcal U_{2\mathrm D}
 =\frac12\vb e^{(0)}:\mathsf A:\vb e^{(0)}
 +\vb e^{(0)}:\mathsf B:\vb{\kappa}
 +\frac12\vb{\kappa}:\mathsf D:\vb{\kappa},
 \label{eq:phys-cont-abd-energy}
\end{equation}
相应的面内力合力 \(\vb N\)（单位长度上的面内力，单位 \(\si{\newton\per\metre}\)）与弯矩合力 \(\vb M\)（单位长度边界上的弯矩，单位 \(\si{\newton}\)）为
\[
 \vb N=\mathsf A:\vb e^{(0)}+\mathsf B:\vb{\kappa},
 \qquad
 \vb M=\mathsf B:\vb e^{(0)}+\mathsf D:\vb{\kappa}.
\]
这就是三维线弹性在 Kirchhoff--Love 运动学下的一般 \(A\!\)-\(B\!\)-\(D\) 母构成。即使 \(\mathsf B\neq0\)，在一个更受限但常用的中间情形中仍可消去膜变量：若 \(\mathsf A\) 在对称面内应变空间上可逆、没有外加面内合力，并且由零膜合力得到的应变场满足 Saint--Venant 相容条件与实际面内边界条件，则 \(\vb N=0\) 给出
\[
 \vb e^{(0)}=-\mathsf A^{-1}:\mathsf B:\vb{\kappa},
 \qquad
 \vb M=\mathsf D_{\rm eff}:\vb{\kappa},
 \qquad
 \mathsf D_{\rm eff}:=\mathsf D-\mathsf B:\mathsf A^{-1}:\mathsf B.
\]
这里 \(\mathsf A^{-1}\) 表示 \(\mathsf A\) 在二维对称应变空间上的逆。这个静力凝聚不是一般的逐点恒等式：对空间非均匀层合板、受约束的面内边界或不相容的候选应变，\(\bm u_t^{(0)}\) 必须由全局膜平衡与相容条件共同求解，不能先局域消去。最干净的纯弯曲闭合仍来自中面对称材料：当材料关于所选中面对称时，\(\mathsf B=0\)，膜与弯曲在线性阶直接解耦；若同时没有面内激励，才可取 \(\vb e^{(0)}=0\) 并进入本章后续的纯弯曲扇区。此时
\begin{equation}
 U[w]
 =\frac12\int_{\mathcal A}
 D_{\alpha\beta\gamma\delta}
 (\partial_\alpha\partial_\beta w)
 (\partial_\gamma\partial_\delta w)\,\dd a,
 \label{eq:phys-cont-general-plate-energy}
\end{equation}
而物理弯矩为
\begin{equation}
 M_{\alpha\beta}
 :=\int_{-h/2}^{h/2}zT_{\alpha\beta}\,\dd z
 =D_{\alpha\beta\gamma\delta}\kappa_{\gamma\delta}
 =-D_{\alpha\beta\gamma\delta}
 \partial_\gamma\partial_\delta w.
 \label{eq:phys-cont-general-plate-moment}
\end{equation}
因此前一章二次型中的 \(\mathsf M_{\mathsf C}[w]\) 对应 \(-\vb M\)（在把二维刚度吸收入 \(\mathsf C\) 后）；场方程和自然边界条件都可以直接用同一个物理弯矩张量表示。

均匀各向同性材料只是\eqnrefs{eq:phys-cont-plane-stress-schur} 的特例。由
\[
 \vb T=2\mu\vb{\varepsilon}+\lambda_{\rm L}(\tr\vb{\varepsilon})\vb I
\]
和 \(T_{zz}=0\) 得到
\[
 \varepsilon_{zz}
 =-\frac{\nu}{1-\nu}\tr\vb{\varepsilon}_t,
\]
继而
\[
 \vb T_t
 =\frac{E}{1-\nu^2}
 \left[(1-\nu)\vb{\varepsilon}_t
 +\nu\tr(\vb{\varepsilon}_t)\vb I_t\right].
\]
于是\eqnrefs{eq:phys-cont-general-plate-moment} 退化为
\begin{equation}
 \vb M
 =D\left[(1-\nu)\vb{\kappa}+\nu\tr(\vb{\kappa})\vb I_t\right],
 \qquad
 D=\frac{Eh^3}{12(1-\nu^2)}.
 \label{eq:phys-cont-kl-moment}
\end{equation}
因为 \(\vb{\kappa}=-\nabla_t\nabla_t w\)，有
\[
 \vb M=-D\left[(1-\nu)\nabla_t\nabla_t w
 +\nu(\nabla_t^2w)\vb I_t\right].
\]
因此前一章的各向同性“曲率应力” \(\mathsf H_\nu[w]\) 与真实弯矩只差物理系数和符号：\(\vb M=-D\mathsf H_\nu[w]\)。

首阶平移动能由面质量密度
\[
 \mu_A(\vb r):=\int_{-h/2}^{h/2}\rho_{\rm m}(\vb r,z)\,\dd z
\]
控制；均匀板有 \(\mu_A=\rho_{\rm m} h\)。若暂时保留由 \(\partial_t\vb u_t=-z\nabla_t\partial_t w\) 产生的转动惯量，则
\begin{equation}
 T[w]
 =\frac12\int_{\mathcal A}
 \left[\mu_A(\partial_tw)^2+I_A|\nabla_t\partial_tw|^2\right]\dd a,
 \qquad
 I_A(\vb r):=\int_{-h/2}^{h/2}\rho_{\rm m}(\vb r,z)z^2\dd z.
 \label{eq:phys-cont-kl-kinetic}
\end{equation}
均匀板满足 \(I_A/\mu_A=h^2/12\)。对波数 \(k\) 的模式，转动惯量相对平移动能的比例因此精确到该运动学阶为 \((kh)^2/12\)。Kirchhoff--Love 首阶理论舍去这一项，同时也舍去同为 \(O((kh)^2)\) 的横向剪切修正。

\begin{proof}[解]
令面内尺度为 $L$、厚度为 $h$，并取
\[
 x_\alpha=L\widehat x_\alpha,\quad z=h\zeta,\quad
 u_z=w_\star\widehat w,\quad
 u_\alpha=\frac{hw_\star}{L}\widehat u_\alpha .
\]
这一标度使面内位移相对横向位移小 $h/L$ 倍。由正文已得的剪切消去\eqnrefs{eq:phys-cont-kl-shear-scaling}、Kirchhoff--Love 运动学\eqnrefs{eq:phys-cont-kl-kinematics}与平面应力 Schur 补\eqnrefs{eq:phys-cont-plane-stress-schur}，三维弹性能首阶即落到中面弯曲扇区。把
\(e_{\alpha\beta}(z)=e^{(0)}_{\alpha\beta}+z\kappa_{\alpha\beta}\)
代入厚度积分，零、一次和二次矩分别产生
\(\mathsf A,\mathsf B,\mathsf D\)。中面对称使奇矩
\(\mathsf B=\int z\mathsf Q\,\dd z\) 严格为零；只有此时纯弯曲扇区与面内伸缩在首阶解耦。
\end{proof}

二维横向载荷必须由三维外力先做厚度约化。对\eqnrefs{eq:phys-cont-3d-elastic-action} 中的体力，首阶横向合力密度为
\[
 f_{\perp}^{(b)}(\vb r,t):=\int_{-h/2}^{h/2}b_z(\vb r,z,t)\,\dd z.
\]
若上下表面还施加外部面牵引 \(\vb t^{+},\vb t^{-}\)，并把它们的 \(z\) 分量都按固定的 \(+z\) 方向计号，则总的单位中面面积横向载荷为
\begin{equation}
 f_\perp=f_{\perp}^{(b)}+t_z^{+}+t_z^{-}.
 \label{eq:phys-cont-transverse-load-reduction}
\end{equation}
侧边牵引不进入这个体载荷，而是在二维极限中成为边界合力与边界弯矩。于是对一般对称板，把\eqnrefs{eq:phys-cont-general-plate-energy} 与首阶动能和 \(f_\perp\) 组合，得到二维作用量
\begin{equation}
 \mathcal S_{2\mathrm D}[w]
 =\int\dd t\left[
 \frac12\int_{\mathcal A}\mu_A(\partial_tw)^2\dd a
 -U[w]+\int_{\mathcal A}f_\perp w\,\dd a
 \right].
 \label{eq:phys-cont-general-plate-action}
\end{equation}
其 Euler--Lagrange 方程为
\begin{equation}
 \mu_A\partial_t^2w
 +\partial_\alpha\partial_\beta
 \left(D_{\alpha\beta\gamma\delta}\partial_\gamma\partial_\delta w\right)
 =\mu_A\partial_t^2w-\partial_\alpha\partial_\beta M_{\alpha\beta}
 =f_\perp.
 \label{eq:phys-cont-general-plate-eom}
\end{equation}
这就是前一章一般四阶算子在薄板中的动力学实现。均匀各向同性时 \(\mathsf D\) 退化为单个 \(D\) 与 \(\nu\)，内部算子中的 \(\nu\) 相消，从而得到
\begin{equation}
 \rho_{\rm m} h\,\partial_t^2w+D\nabla_t^4w=f_\perp,\qquad
 \nabla_t^4=(\nabla_t^2)^2.
 \label{eq:phys-cont-kl-eom}
\end{equation}
静力、自由振动和受迫振动只是\eqnrefs{eq:phys-cont-general-plate-eom} 的不同时间与载荷选择，并不是三套独立理论。

边界量同样来自作用量。沿光滑边界段记外法向和切向为 \(\vb n,\vb\tau\)。对物理弯矩 \(\vb M\) 定义 Kirchhoff 剪力和有效剪力
\begin{equation}
 \bm S_{\rm K}:=\nabla_t\cdot\vb M,
 \qquad
 V_n:=S_{{\rm K},n}+\partial_sM_{n\tau}.
 \label{eq:phys-cont-kl-effective-shear}
\end{equation}
边界虚功只包含 \((\partial_nw,M_{nn})\) 与 \((w,V_n)\) 两对共轭量，因此经典边界条件为
\begin{align*}
 \mathrm C:&\quad w=0,\qquad \partial_nw=0,\\
 \mathrm S:&\quad w=0,\qquad M_{nn}=0,\\
 \mathrm F:&\quad M_{nn}=0,\qquad V_n=0.
\end{align*}
在分段光滑边界的角点 \(P_j\)，沿边界的切向分部积分还留下
\[
 \bigl(M_{n\tau}^{(j,-)}-M_{n\tau}^{(j,+)}\bigr)\delta w(P_j).
\]
若没有集中角力，该跳跃必须为零。于是场方程、自由边界、角点平衡都来自同一个二维作用量；这也解释了为什么这些边界条件恰好与前一章的自伴 Green 边界型相容。

\begin{proof}[解]
令
\[
 N_{\alpha\beta}
 :=D_{\alpha\beta\gamma\delta}\partial_\gamma\partial_\delta w.
\]
由\eqnrefs{eq:phys-cont-general-plate-energy}，
\[
 \delta U
 =\int_{\mathcal A}N_{\alpha\beta}\partial_\alpha\partial_\beta\delta w\,\dd a.
\]
按前一章的一般 Green 恒等式两次分部积分，
\begin{align*}
 \delta U
 ={}&\int_{\mathcal A}
 \partial_\alpha\partial_\beta N_{\alpha\beta}\,\delta w\,\dd a\\
 &+\int_{\partial\mathcal A}
 \left[N_{nn}\partial_n\delta w
 -\left(n_\beta\partial_\alpha N_{\alpha\beta}
 +\partial_sN_{n\tau}\right)\delta w\right]\dd s
 +\text{corner terms}.
\end{align*}
时间动能满足
\[
 \delta T
 =-\int_{\mathcal A}\mu_A\partial_t^2w\,\delta w\,\dd a
\]
（时间端点变分取零），外载虚功为 \(\int f_\perp\delta w\dd a\)。因此 Hamilton 原理的体积分系数给出
\[
 \mu_A\partial_t^2w
 +\partial_\alpha\partial_\beta N_{\alpha\beta}=f_\perp,
\]
即正文\eqnrefs{eq:phys-cont-general-plate-eom}。

由一般弯矩定义\eqnrefs{eq:phys-cont-general-plate-moment}，\(N=-\vb M\)，故
\[
 N_{nn}=-M_{nn},
 \qquad
 n\cdot(\nabla\cdot N)+\partial_sN_{n\tau}=-V_n.
\]
边界项只整体差一个符号约定，不改变自由条件 \(M_{nn}=V_n=0\)。把
\[
 \vb M=-D\left[(1-\nu)\nabla_t\nabla_t w
 +\nu(\nabla_t^2w)\vb I_t\right]
\]
代入体方程，有
\begin{align*}
 -\partial_\alpha\partial_\beta M_{\alpha\beta}
 &=D\left[(1-\nu)\partial_\alpha\partial_\beta
 \partial_\alpha\partial_\beta w
 +\nu\partial_\alpha\partial_\alpha\nabla_t^2w\right]\\
 &=D\bigl[(1-\nu)+\nu\bigr]\nabla_t^4w
 =D\nabla_t^4w,
\end{align*}
于是恢复\eqnrefs{eq:phys-cont-kl-eom}。
\end{proof}

薄参数 \(h/L\) 只控制降维，并不自动保证所有其他近似。对典型波数 \(k\)，Kirchhoff--Love 长波条件是 \(kh\ll1\)。均匀板的转动惯量比已经由\eqnrefs{eq:phys-cont-kl-kinetic} 给出
\[
 \frac{T_{\rm rotary}}{T_{\rm transverse}}
 =\frac{h^2k^2}{12},
\]
横向剪切导致的 Mindlin--Reissner 修正也从 \(O((kh)^2)\) 进入，因此 \(kh\) 是频域中比“几何上看起来很薄”更直接的控制参数。

线性几何还要求另一组独立条件。典型斜率是 \(w_\star/L\)，von K\'arm\'an 中面应变约为 \((w_\star/L)^2\)，而线性弯曲应变约为 \(hw_\star/L^2\)，两者比值为
\begin{equation}
 \frac{\varepsilon_{\rm geom}}{\varepsilon_{\rm bend}}
 \sim\frac{w_\star}{h}.
 \label{eq:phys-cont-kl-geometric-control}
\end{equation}
因此 \(w_\star\gtrsim h\) 时，即使 \(w_\star/L\ll1\)，由几何非线性诱导的中面拉伸也会与线性弯曲同阶，不能把它与线性层合板的 \(\mathsf B\) 耦合混为一谈；此时应转向 von K\'arm\'an 或更完整的壳/板理论。

最后，二维外层展开在距边缘和角点 \(O(h)\) 的三维边界层内通常不逐点有效；固支边的局部剪切、角点应力集中以及厚度高阶模会在那里恢复。若 \(kh\not\ll1\)，应采用含剪切和转动惯量的 Mindlin--Reissner 理论或直接回到三维弹性；若厚度方向缺乏中面对称性，则\eqnrefs{eq:phys-cont-general-plate-rigidity} 中一般有 \(\mathsf B\neq0\)。只有在零膜合力且相容与面内边界条件允许时，才可用前述 \(\mathsf D_{\rm eff}\) 做静力凝聚；一般情况下仍必须保留\eqnrefs{eq:phys-cont-abd-energy} 的膜--弯曲耦合。若进入塑性、粘弹或大应变区，则线性构成与\eqnrefs{eq:phys-cont-3d-elastic-action} 本身都需要替换。以下环板和长方形只是在这些条件满足后，再进一步施加均匀各向同性材料与高对称几何得到的特例。
